\section{The inactive vertex distribution}
\label{sec:marked}
This entire section assumes fixed $p\ge2$.
Let $R_n$ be the inactive-vertex count from Section~\ref{sec:exact}, let
$a\equiv n\pmod q$, and let $Q_{\nu,a}$ be the distribution of
$\Pois(\nu)$ conditioned on the residue class $a$ modulo $q$.
For probability measures on this lattice we use
\[
 \TV(P,Q)=\frac12\sum_r|P(r)-Q(r)|.
\]
The exact Poisson identity shows that the mass of $R_n$ is proportional to
\begin{equation}\label{eq:markedweight}
 \frac{\lambda^r}{r!}e^{\Delta_n(r)},
 \qquad 0\le r\le n,\quad r\equiv a\pmod q.
\end{equation}
All constants $A,b,\gamma_1,\gamma_2,\gamma_3$ retain the meanings in
\eqref{eq:firstuniv} and \eqref{eq:D2}.

\subsection{The two Poisson references}
\begin{theorem}\label{thm:markedmain}
There is a unique positive solution $\mu$ to
\begin{equation}\label{eq:mu}
 \log(\mu/\lambda)=-\frac{2A\mu+b}{n}.
\end{equation}
In terms of the principal positive branch of the Lambert function,
\begin{equation}\label{eq:muW}
 \mu=\frac{n}{2A}W\left(\frac{2A\lambda}{n}e^{-b/n}\right)
 =\lambda-\frac{2A\lambda^2+b\lambda}{n}
                   +O_p(\lambda^3/n^2).
\end{equation}
The mean and variance satisfy
\begin{align}
 \E R_n&=\lambda-\frac{2A\lambda^2+b\lambda}{n}
                      +O_p(\lambda^3/n^2),\label{eq:mean}\\
 \Var R_n&=\lambda-\frac{4A\lambda^2+b\lambda}{n}
                      +O_p(\lambda^3/n^2).\label{eq:variance}
\end{align}
Moreover,
\begin{align}
 \TV(\mathcal L(R_n),Q_{\lambda,a})
       &\sim A\sqrt{2/\pi}\,\frac{\lambda^{3/2}}{n},\label{eq:TVlambda}\\
 \TV(\mathcal L(R_n),Q_{\mu,a})
       &\sim A\sqrt{2/(\pi e)}\,\frac{\mu}{n}.\label{eq:TVmu}
\end{align}
The estimates and equivalents hold uniformly over the residue classes.
\end{theorem}

Uniqueness in \eqref{eq:mu} follows because
$\log\mu+2A\mu/n$ is strictly increasing from $-\infty$ to $+\infty$.
Equation \eqref{eq:muW} follows from $W(z)e^{W(z)}=z$;
its expansion uses $\lambda/n\to0$.
Since $\lambda^2/n\to0$, the remainder in \eqref{eq:mean} and
\eqref{eq:variance} is smaller even than the displayed $\lambda/n$ terms.
In particular,
\begin{equation}\label{eq:fano}
 \frac{\Var R_n}{\E R_n}
       =1-\frac{2A\lambda}{n}+O_p(\lambda^2/n^2).
\end{equation}

The reference in both total-variation statements is conditioned.
An unconditioned Poisson distribution gives mass approaching $1-1/q$ to
the complement of the support of $R_n$, so its total variation from
$R_n$ is at least $1-1/q+o(1)$. A Gaussian limit after scaling and a
small total-variation distance on the original integer lattice are
therefore different assertions.

\subsection{A density expansion with weighted remainder}
Define the monic second Poisson--Charlier polynomial by
\begin{equation}\label{eq:charlier}
 \mathcal C_2(r;\mu)=(r-\mu)^2-r.
\end{equation}
This normalization is $\mu^2$ times the degree-two Charlier polynomial
in the generating-function convention of DLMF 18.23.5~\cite{dlmf}.
Let
\begin{equation}\label{eq:K2}
 K_2(r;\mu)=D_2(r)+\frac{A^2}{2}\mathcal C_2(r;\mu)^2.
\end{equation}
Write $\E_\mu$ for ordinary Poisson expectation and $\E_Q$ for expectation
under $Q_{\mu,a}$.

\begin{theorem}\label{thm:weighted}
On the support of $Q_{\mu,a}$ define
$g_n(r)=\Pp(R_n=r)/Q_{\mu,a}(r)$, taking $g_n(r)=0$ for $r>n$.
Then
\begin{align}\label{eq:secondensity}
 g_n(r)={}&1-\frac{A}{n}\mathcal C_2(r;\mu)\nonumber\\
 &+\frac{K_2(r;\mu)-\E_\mu K_2(X;\mu)}{n^2}
   +\epsilon_n(r),
\end{align}
where, for $s=0,1,2$,
\begin{equation}\label{eq:weightedrem}
 \E_Q[|X-\mu|^s|\epsilon_n(X)|]
                       =O_p(\mu^{4+s/2}/n^3).
\end{equation}
Consequently,
\begin{equation}\label{eq:firstdensity}
 \E_Q\left|g_n(X)-1+\frac A n\mathcal C_2(X;\mu)\right|
                          =O_p(\mu^{5/2}/n^2).
\end{equation}
For each fixed $L>0$, uniformly at allowed points with
$|r-\mu|\le L\sqrt\mu$,
\begin{equation}\label{eq:localdensity}
 g_n(r)=1-\frac A n\mathcal C_2(r;\mu)
                         +O_{p,L}(\mu^{5/2}/n^2).
\end{equation}
\end{theorem}

\begin{proof}
Write $\beta_0=\log(\mu/\lambda)=-(2A\mu+b)/n$.
Changing the reference Poisson parameter subtracts $\beta_0r$ from the
exponent in \eqref{eq:markedweight}, up to a normalizing constant.
The identity $q/2+b=A$ gives the exact cancellation
\begin{equation}\label{eq:cancel}
 \frac{D_1(r)}n-\beta_0r
       =\frac{A\mu^2}{n}-\frac A n\mathcal C_2(r;\mu).
\end{equation}
The constant disappears after normalization. On $r\le2\mu$, logarithmic
Taylor expansion gives
\[
 \Delta_n(r)=\frac{D_1(r)}n+\frac{D_2(r)}{n^2}
                                  +O_p(r^4/n^3).
\]
Since $\mu/\lambda\to1$ and $\mu^2/n\to0$, the residual exponent is
uniformly $o(1)$ on this interval:
\begin{equation}\label{eq:residualV}
 V(r)=-\frac A n\mathcal C_2(r;\mu)
              +\frac{D_2(r)}{n^2}+O_p(r^4/n^3).
\end{equation}
Its exponential is
$1-A\mathcal C_2/n+K_2/n^2$ plus a remainder.

We give the weighted estimates explicitly. For $Y=X-\mu$, the centered
Poisson cumulant generating function shows that
$\E_\mu Y^{2h}=O_h(\mu^h)$ for $\mu\ge1$.
Cauchy--Schwarz and these even moments imply, for $s=0,1,2$,
\begin{align*}
 \E_\mu[|Y|^s|\mathcal C_2|^3]&=O_p(\mu^{3+s/2}),\\
 \E_\mu[|Y|^s|\mathcal C_2D_2|]&=O_p(\mu^{4+s/2}),\\
 \E_\mu[|Y|^sD_2^2]&=O_p(\mu^{6+s/2}),\\
 \E_\mu[|Y|^sX^4]&=O_p(\mu^{4+s/2}).
\end{align*}
For example $\mathcal C_2=Y^2-Y-\mu$ and $X=\mu+Y$ reduce each assertion
to a fixed collection of central moments. Conditioning changes upper
bounds only by a fixed factor because
$\Pp_\mu(X\equiv a\pmod q)\to1/q$.
The four displayed bounds control respectively the cubic exponential
term, the cross term, the square of $D_2$, and the exponent remainder.
The $D_2^2$ contribution carries $n^{-4}$ and is absorbed into
$O(\mu^{4+s/2}/n^3)$ because $\mu^2/n\to0$.

Tails beyond $2\mu$ are exponentially small after any fixed polynomial
weight. For the exact law, this follows from the original
$e^{\Delta_n(r)}\le1$ domination under $\Pois(\lambda)$, the normalizer
tending to one, and $2\mu/\lambda\to2$.
For the reference it is an ordinary Chernoff bound.
Thus no global pointwise estimate on the changed-reference density is
being assumed.

It remains to check normalization. The exact identities
$\E_\mu\mathcal C_2=0$ and
$\E_\mu\mathcal C_2^2=2\mu^2$ give
$\E_\mu K_2=O_p(\mu^3)$.
The fixed-polynomial residue correction is exponentially small by
\eqref{eq:filter}. Dividing by the normalizer subtracts
$\E_\mu K_2/n^2$ at second order. Its cross term with
$\mathcal C_2/n$ has weighted size
$O(\mu^{4+s/2}/n^3)$, and the normalizer-square terms have size
$O(\mu^{6+s/2}/n^4)$, again absorbed by the same bound.
This proves \eqref{eq:secondensity}--\eqref{eq:weightedrem}.

Finally, $D_2-\E_\mu D_2$ has $L^1$ norm $O_p(\mu^{5/2})$,
while $\mathcal C_2^2-\E_\mu\mathcal C_2^2$ has $L^1$ norm
$O(\mu^2)$. These follow by expanding in $Y$ and using its central moments.
They prove \eqref{eq:firstdensity}; its third-order remainder is smaller
because $\mu^{3/2}/n\to0$.
On $|r-\mu|\le L\sqrt\mu$, the corresponding pointwise centered bounds
hold, and the normalized third-order remainder is
$O_{p,L}(\mu^4/n^3)=o(\mu^{5/2}/n^2)$.
This proves \eqref{eq:localdensity}.
\end{proof}

\subsection{Moments and their second correction}
The polynomial identities needed to extract the moments are
\begin{align}
 \E_\mu\mathcal C_2&=0,&
 \E_\mu[(X-\mu)\mathcal C_2]&=0,\label{eq:charlierorth}\\
 \E_\mu[((X-\mu)^2-\mu)\mathcal C_2]&=2\mu^2,&
 \E_\mu[(X-\mu)\mathcal C_2^2]&=4\mu^2,\nonumber\\
 \E_\mu[((X-\mu)^2-\mu)\mathcal C_2^2]
        &=8\mu^3+8\mu^2.&&\nonumber
\end{align}
They can be verified directly by the factorial moment rule.
More generally, Poisson summation by parts gives for any polynomial $f$
\begin{align}
 \Cov_\mu(X,f(X))&=\mu\E_\mu[f(X+1)-f(X)],\label{eq:parts1}\\
 \Cov_\mu((X-\mu)^2,f(X))
  &=\mu\E_\mu[f(X+1)-f(X)]\nonumber\\
  &\quad+\mu^2\E_\mu[f(X+2)-2f(X+1)+f(X)].\label{eq:parts2}
\end{align}
For example, $\E[Xf(X)]=\mu\E f(X+1)$ gives the first identity;
using $\E[(X)_2f(X)]=\mu^2\E f(X+2)$ gives the second.

Multiply \eqref{eq:secondensity} by $X-\mu$ or $(X-\mu)^2$.
The second-order polynomial terms are $O_p(\mu^3/n^2)$ by
\eqref{eq:charlierorth}--\eqref{eq:parts2}.
The weighted errors are $O_p(\mu^{9/2}/n^3)$ and
$O_p(\mu^5/n^3)$, respectively, both $o(\mu^3/n^2)$.
The small centered mean squared is absorbed in the variance error.
Consequently
\begin{equation}\label{eq:meanmu}
 \E R_n=\mu+O_p(\mu^3/n^2),\qquad
 \Var R_n=\mu-\frac{2A\mu^2}{n}+O_p(\mu^3/n^2).
\end{equation}
Substitution of \eqref{eq:muW} proves \eqref{eq:mean} and
\eqref{eq:variance}.

A useful second-order form is obtained by setting $\theta=z\,d/dz$ and
\begin{align}\label{eq:L2}
 L_2(z)={}&(\gamma_3+2A^2)z^3
 +(3\gamma_3-\gamma_2+A^2+2Ab)z^2\nonumber\\
 &+(\gamma_3-\gamma_2-\gamma_1+b^2/2)z.
\end{align}
Then the same weighted calculation, retaining its polynomial terms, gives
\begin{align}
 \E R_n={}&\lambda-\frac{2A\lambda^2+b\lambda}{n}
            +\frac{(\theta L_2)(\lambda)}{n^2}
            +O_p(\lambda^{9/2}/n^3),\label{eq:mean2}\\
 \Var R_n={}&\lambda-\frac{4A\lambda^2+b\lambda}{n}
            +\frac{(\theta^2L_2)(\lambda)}{n^2}
            +O_p(\lambda^5/n^3).\label{eq:var2}
\end{align}
These are conservative weighted remainder bounds.
An independent way to verify the coefficient is to write
$D_1=-A(X)_2-bX$ and use
\[
 L_2=\E D_2+\tfrac12\Var(D_1),\quad
 \Var((X)_2)=4\lambda^3+2\lambda^2,\quad
 \Cov((X)_2,X)=2\lambda^2.
\]
It yields exactly \eqref{eq:L2}. In the two highlighted cases,
\begin{align*}
 p=2:\quad &L_2(z)=12z^3+\frac{44}{3}z^2+\frac43z,\\
 p=3:\quad &L_2(z)=\frac{4967}{96}z^3
                      +\frac{4515}{64}z^2+\frac{1161}{128}z.
\end{align*}

\subsection{Sharp total variation constants}
\begin{proof}[Proof of \eqref{eq:TVlambda} and \eqref{eq:TVmu}]
Expansion at the original reference and division by its normalizer give
\begin{equation}\label{eq:originaldensity}
 \frac{\Pp(R_n=r)}{Q_{\lambda,a}(r)}
 =1+\frac{D_1(r)-\E_{Q_{\lambda,a}}D_1}{n}+\rho_n(r),
 \qquad \E_{Q_{\lambda,a}}|\rho_n|=O_p(\lambda^4/n^2).
\end{equation}
The same local expansion and global tail proof as in Section~\ref{sec:subcritical}
justifies this $L^1$ statement, including the cutoff.
Under $Q_{\lambda,a}$, write $Z_\lambda=(X-\lambda)/\sqrt\lambda$.
Then $Z_\lambda$ converges to a standard normal $Z$, uniformly over residues,
and its fixed even moments are bounded.
To see convergence, apply the residue filter to the characteristic function
at $t/\sqrt\lambda$. The zero root gives the ordinary Poisson central
limit, and every nonzero root decays exponentially for bounded $t$.
Moment bounds follow from the fixed-polynomial filter.

Expansion of $D_1$ about $\lambda$ now gives
\[
 \frac{D_1(X)-\E_{Q_{\lambda,a}}D_1}{\lambda^{3/2}}
                    \longrightarrow -2AZ
\]
in distribution and in absolute first moment, by uniform integrability.
Since $\lambda^{5/2}/n\to0$ for $p\ge2$, the error in
\eqref{eq:originaldensity} is negligible relative to
$\lambda^{3/2}/n$. Thus
\[
 \TV(\mathcal L(R_n),Q_{\lambda,a})
     \sim\frac{2A\lambda^{3/2}}{2n}\E|Z|
     =A\sqrt{2/\pi}\,\frac{\lambda^{3/2}}n.
\]

For the tilted reference, \eqref{eq:firstdensity} gives
\[
 \TV(\mathcal L(R_n),Q_{\mu,a})
 =\frac{A}{2n}\E_{Q_{\mu,a}}|\mathcal C_2(X;\mu)|
                                    +O_p(\mu^{5/2}/n^2).
\]
The same conditional central limit and uniform integrability imply
$\mathcal C_2(X;\mu)/\mu\to Z^2-1$ in distribution and absolute first
moment. If $\phi$ is the standard normal density, integration by parts
on $[0,1]$ gives $\E|Z^2-1|=4\phi(1)$.
Since $\mu^{3/2}/n\to0$, the error is negligible, and the leading constant is
$2A\phi(1)=A\sqrt{2/(\pi e)}$.
\end{proof}

\subsection{All fixed marked orders and graph consequences}
The finite polynomials $U_j$ in \eqref{eq:U} give an algebraic logarithm
by the recurrence
\begin{equation}\label{eq:Lrec}
 L_j(z)=U_j(z)-\frac1j\sum_{k=1}^{j-1}kL_k(z)U_{j-k}(z).
\end{equation}
Equivalently, as a formal series in $w$,
\[
 \log\left(1+\sum_{j\ge1}U_j(z)w^j\right)
                      =\sum_{j\ge1}L_j(z)w^j.
\]
The formal coefficient of the $h$th cumulant of $R_n$ is then
\begin{equation}\label{eq:cumulants}
 \lambda+\sum_{j\ge1}\frac{(\theta^hL_j)(\lambda)}{n^j}.
\end{equation}
This is an effective all-fixed-order generator, with the following precise
remainder interpretation. Multiply the finite density expansion through
$J$ by each required fixed power $r^d$, apply the polynomial filter, and
form the finite quotients and cumulant combinations. The unnormalized
weighted error is bounded by
\[
 O_{p,J,d}(\lambda^{2J+2+d}/n^{J+1}).
\]
For a desired fixed algebraic order, take $J$ sufficiently large and then
collect powers after setting $\lambda=c_p n^{1/q}$.
This procedure justifies \eqref{eq:cumulants} to that order without
assuming any general degree cancellation in $L_j$, and without
differentiating an uncontrolled asymptotic remainder.

For clarity, an exact marked generating function can be written before
any approximation:
\[
 \E e^{sR_n}=
 \frac{\sum_{r\equiv a\ (q),\,r\le n}(\lambda e^s)^r e^{\Delta_n(r)}/r!}
      {\sum_{r\equiv a\ (q),\,r\le n}\lambda^r e^{\Delta_n(r)}/r!}.
\]
Its finite derivatives explain the operator $\theta$ in the coefficient
rule. The error control is supplied by the weighted argument above.

Finally, if $K_n$ and $V_n$ are the numbers of active and total vertices,
\eqref{eq:rdef} gives the exact relations
\[
 \E K_n=\frac{n-\E R_n}{q},\quad
 \Var K_n=\frac{\Var R_n}{q^2},\quad
 \E V_n=\frac{n+p\E R_n}{q},\quad
 \Var V_n=\frac{p^2\Var R_n}{q^2}.
\]
The two vertex counts are affinely determined by the same defect.
Their centered Gaussian limits follow from the conditional Poisson limit
and either total-variation approximation. None of the statements in
this section is asserted for $p=1$.
